\chapter{Topological Sort and DAGs}

In Chapter~3, depth-first search timestamps showed us how to detect cycles:
a directed graph has a cycle if and only if DFS encounters a back edge.
When no directed cycles exist, the graph is a directed acyclic graph -- a DAG --
the core model of dependency. Every time you compile code, resolve package
dependencies, or schedule tasks, you are walking a DAG. The central operation on
a DAG is to linearize its partial order into a sequence respecting every prerequisite:
a topological order. Once vertices are topologically ordered, problems that are
NP-hard on general graphs (such as longest paths) collapse into linear-time dynamic
programs. This chapter builds topological orderings via DFS finish times and Kahn's
indegree-peeling algorithm, proves both methods correct, and applies the order to
dynamic programming.

\chapterline

\section{What a DAG Is, and Why Acyclicity Unlocks Linear Order}

\begin{definition}[Directed acyclic graph]\label{ch04:dag-def}
A directed graph $\Gph=(V,E)$ is a \emph{directed acyclic graph} (DAG) if it contains no directed
cycle: there is no sequence of vertices $v_0,v_1,\dots,v_k=v_0$ with $k\ge 1$ and every
$(v_{i-1},v_i)\in E$.
\end{definition}

The prohibition is on \emph{directed} cycles. An undirected triangle is fine to orient as
$a\to b\to c$, $a\to c$; nothing closes a loop. What kills a DAG is a back-pointing edge that lets
you return to where you started.

\begin{intuitionbox}
Read each edge $u\to v$ as ``$u$ must come before $v$,'' or ``$v$ depends on $u$.'' A directed cycle
$u\to\cdots\to u$ then says ``$u$ must come strictly before itself,'' which is impossible to satisfy.
So acyclicity is precisely the condition that the constraints can \emph{all} be met at once. A DAG is
a consistent set of ``before'' rules; topological sort is the act of obeying them.
\end{intuitionbox}

\begin{definition}[Topological order]\label{ch04:topo-def}
A \emph{topological order} of a digraph $\Gph=(V,E)$ is a linear order $\prec$ on $V$ (equivalently,
a bijection $\pi\colon V\to\set{1,\dots,n}$) such that for every edge $(u,v)\in E$ we have
$\pi(u)<\pi(v)$: every edge points forward in the order.
\end{definition}

If we lay the vertices on a horizontal line in topological order, \emph{every} arrow points to the
right. That picture is worth memorizing.

\begin{center}
\begin{tikzpicture}[>=Latex, node distance=14mm,
  v/.style={circle,draw=nodeblue,fill=blue!6,thick,minimum size=8mm,inner sep=0pt},
  every edge/.style={draw=edgegray,thick,->}]
  \node[v] (a) {$a$};
  \node[v,right=of a] (b) {$b$};
  \node[v,right=of b] (c) {$c$};
  \node[v,right=of c] (d) {$d$};
  \node[v,right=of d] (e) {$e$};
  \draw (a) edge (b);
  \draw (b) edge (c);
  \draw (c) edge (d);
  \draw (d) edge (e);
  \draw (a) edge[bend left=35] (c);
  \draw (b) edge[bend right=40] (d);
  \draw (a) edge[bend left=45] (e);
\end{tikzpicture}
\end{center}

\begin{focusbox}
Two facts drive everything in this chapter:
\begin{enumerate}
\item A digraph has a topological order \emph{if and only if} it is a DAG.
\item Both standard algorithms (DFS finish times, Kahn indegree peeling) run in $\Oh(n+m)$ and
either produce an order or expose a cycle.
\end{enumerate}
\end{focusbox}

\begin{theorem}[Existence]\label{ch04:existence}
A finite digraph $\Gph$ has a topological order if and only if it is acyclic.
\end{theorem}

\begin{proof}
($\Rightarrow$) Suppose $\pi$ is a topological order and, for contradiction, $\Gph$ has a cycle
$v_0\to v_1\to\cdots\to v_k=v_0$. Each edge forces $\pi(v_{i-1})<\pi(v_i)$, so chaining gives
$\pi(v_0)<\pi(v_1)<\cdots<\pi(v_k)=\pi(v_0)$, i.e.\ $\pi(v_0)<\pi(v_0)$, impossible.

($\Leftarrow$) Suppose $\Gph$ is acyclic. We build the order greedily, which also proves the next
lemma. Pick any vertex with no incoming edge (Lemma~\ref{ch04:source-exists} below guarantees one
exists), assign it the smallest unused position, delete it, and repeat on the rest. Deleting a
vertex cannot create a cycle, so every subgraph is again a DAG and again has a source. After $n$
steps every vertex has a position, and any surviving edge $(u,v)$ had $u$ removed before $v$ (we only
remove a vertex once it has no in-edges from undeleted vertices), so $\pi(u)<\pi(v)$.
\end{proof}

\begin{lemma}[Every finite DAG has a source and a sink]\label{ch04:source-exists}
A nonempty finite DAG contains at least one \emph{source} (a vertex of indegree $0$) and at least one
\emph{sink} (a vertex of outdegree $0$).
\end{lemma}

\begin{proof}
We show a source exists; the sink argument is symmetric (reverse all edges). Suppose every vertex had
an incoming edge. Start at any $v_0$, follow an in-edge backward to a predecessor $v_1$, then a
predecessor $v_2$ of $v_1$, and so on. Each step picks a predecessor, which exists by assumption.
After $n$ steps we have visited $n+1$ vertices $v_0,\dots,v_n$, so by pigeonhole two coincide, say
$v_i=v_j$ with $i<j$. The backward walk $v_i\leftarrow\cdots\leftarrow v_j$ reversed is a forward
directed cycle $v_j\to\cdots\to v_i=v_j$ -- contradicting acyclicity. Hence some vertex has indegree
$0$.
\end{proof}

\begin{whybox}
Lemma~\ref{ch04:source-exists} is the load-bearing fact. ``A DAG always has somewhere to start'' is
why both algorithms make progress: peeling needs a source to remove; DFS needs a vertex that finishes
with all successors already processed (a ``local sink''). Without acyclicity neither holds, and that
failure is exactly how the algorithms detect a cycle.
\end{whybox}

\begin{historybox}
Topological ordering was formalized by A.~B.~Kahn in 1962 in a paper on scheduling project networks
(the PERT/CPM tradition of operations research). The depth-first formulation is folklore but is most
cleanly stated through the finish-time machinery that Tarjan and Hopcroft popularized in the early
1970s, the same timestamp framework that powers strongly connected components and biconnectivity. The
two views -- ``peel from the front'' and ``recurse and reverse'' -- are duals that we will meet again
and again.
\end{historybox}

\section{Method 1: Kahn's Algorithm (Indegree Peeling)}

The constructive proof of Theorem~\ref{ch04:existence} \emph{is} an algorithm. Repeatedly: take a
vertex with indegree $0$, output it, and remove it (which decrements the indegree of its
out-neighbors). The only engineering question is how to find zero-indegree vertices fast. The trick:
maintain a running indegree count for every vertex and keep a queue of those currently at $0$.

\begin{intuitionbox}
Think of the DAG as a stack of dependencies and you are clearing a to-do list. A task is
\emph{ready} when all its prerequisites are done -- i.e.\ its current indegree is $0$. Do a ready
task, which may make some other tasks ready, and add those to the ready set. The order in which tasks
leave the ready set is a valid topological order. If at some point the list is nonempty but
\emph{nothing} is ready, the leftover tasks form a cycle of mutual prerequisites.
\end{intuitionbox}

\begin{algorithm}[H]
\DontPrintSemicolon
\KwIn{Digraph $\Gph=(V,E)$ as adjacency lists, $n=\abs V$, $m=\abs E$.}
\KwOut{A topological order $L$, or report ``cycle.''}
\ForEach{$v\in V$}{$\mathrm{indeg}[v]\gets 0$}
\ForEach{edge $(u,v)\in E$}{$\mathrm{indeg}[v]\gets \mathrm{indeg}[v]+1$}
$Q\gets$ empty queue\;
\ForEach{$v\in V$ with $\mathrm{indeg}[v]=0$}{enqueue $v$ into $Q$}
$L\gets$ empty list\;
\While{$Q$ is not empty}{
  $u\gets$ dequeue($Q$)\;
  append $u$ to $L$\;
  \ForEach{$w$ with $(u,w)\in E$}{
    $\mathrm{indeg}[w]\gets \mathrm{indeg}[w]-1$\;
    \If{$\mathrm{indeg}[w]=0$}{enqueue $w$ into $Q$}
  }
}
\eIf{$\abs L = n$}{\Return $L$}{\Return ``cycle'' (the $n-\abs L$ unoutput vertices lie on/after a cycle)}
\caption{\textsc{KahnTopoSort}$(\Gph)$}\label{ch04:kahn-algo}
\end{algorithm}

\subsection*{Worked trace}

Take a concrete DAG of seven tasks. Read edges as ``must precede.''

\begin{center}
\begin{tikzpicture}[>=Latex, node distance=16mm,
  v/.style={circle,draw=nodeblue,fill=blue!6,thick,minimum size=8.5mm,inner sep=0pt},
  every edge/.style={draw=edgegray,thick,->}]
  \node[v] (a) {$1$};
  \node[v,below right=10mm and 16mm of a] (b) {$2$};
  \node[v,above right=10mm and 16mm of a] (c) {$3$};
  \node[v,right=34mm of a] (d) {$4$};
  \node[v,right=of d] (e) {$5$};
  \node[v,above right=10mm and 16mm of e] (f) {$6$};
  \node[v,below right=10mm and 16mm of e] (g) {$7$};
  \draw (a) edge (b);
  \draw (a) edge (c);
  \draw (b) edge (d);
  \draw (c) edge (d);
  \draw (c) edge[bend left=12] (f);
  \draw (d) edge (e);
  \draw (e) edge (f);
  \draw (e) edge (g);
  \draw (d) edge[bend right=20] (g);
\end{tikzpicture}
\end{center}

Initial indegrees: $1{:}0,\ 2{:}1,\ 3{:}1,\ 4{:}2,\ 5{:}1,\ 6{:}2,\ 7{:}2$. The only source is $1$,
so $Q=[1]$. We process FIFO; ties break by smaller vertex number for determinism.

\begin{center}
\renewcommand{\arraystretch}{1.2}
\begin{tabular}{r l l l}
\toprule
step & dequeue & decrements (new indeg) & queue after \\
\midrule
1 & $1$ & $2{:}0$, $3{:}0$ & $[2,3]$ \\
2 & $2$ & $4{:}1$ & $[3]$ \\
3 & $3$ & $4{:}0$, $6{:}1$ & $[4]$ \\
4 & $4$ & $5{:}0$, $7{:}1$ & $[5]$ \\
5 & $5$ & $6{:}0$, $7{:}0$ & $[6,7]$ \\
6 & $6$ & --- & $[7]$ \\
7 & $7$ & --- & $[\,]$ \\
\bottomrule
\end{tabular}
\end{center}

Output order: $L=(1,2,3,4,5,6,7)$. All seven vertices appear, so no cycle. Sanity check: every edge
points forward, e.g.\ $3\to 6$ has $3$ before $6$, and $4\to 7$ has $4$ before $7$.
The figure below contrasts how Kahn's indegree peeling wave advances through this DAG (left) against
how a cyclic dependency traps its vertices with positive residual indegree so the queue runs dry early (right).

\begin{center}
\begin{tikzpicture}[
    >=Stealth,
    vdone/.style={circle, draw=edgegray, fill=gray!12, thick, minimum size=7.2mm, inner sep=0pt, font=\small},
    vq/.style={circle, draw=targetgreen!85!black, fill=targetgreen!18, very thick, minimum size=7.2mm, inner sep=0pt, font=\small\bfseries},
    vwait/.style={circle, draw=sourceblue, fill=sourceblue!8, thick, minimum size=7.2mm, inner sep=0pt, font=\small},
    vstuck/.style={circle, draw=trapred!90!black, fill=trapred!15, very thick, minimum size=7.2mm, inner sep=0pt, font=\small\bfseries},
    peeled/.style={draw=gray!55, dashed, thick, ->},
    active/.style={draw=edgegray, thick, ->},
    cyc/.style={draw=trapred!85!black, very thick, ->},
    paneltitle/.style={font=\footnotesize\scshape, text=gray!85!black},
    badge/.style={font=\tiny\ttfamily, inner sep=1pt}
]
  % Left panel: Kahn peeling snapshot after step 1 (1 dequeued, 2 and 3 in Q)
  \begin{scope}[xshift=0cm]
    \node[paneltitle] at (3.1, 2.35) {(a) After peeling $1$: $Q=[2,3]$ unlocked};

    \node[vdone] (k1) at (0.0, 0.6) {$1$};
    \node[vq]    (k2) at (1.5, -0.3) {$2$};
    \node[vq]    (k3) at (1.5, 1.5) {$3$};
    \node[vwait] (k4) at (3.1, 0.6) {$4$};
    \node[vwait] (k5) at (4.6, 0.6) {$5$};
    \node[vwait] (k6) at (6.1, 1.5) {$6$};
    \node[vwait] (k7) at (6.1, -0.3) {$7$};

    \node[badge, text=gray!70!black, above=2pt of k1] {out};
    \node[badge, text=targetgreen!85!black, below=2pt of k2] {indeg=0};
    \node[badge, text=targetgreen!85!black, above=2pt of k3] {indeg=0};
    \node[badge, text=sourceblue!85!black, above=2pt of k4] {indeg=2};
    \node[badge, text=sourceblue!85!black, above=2pt of k5] {indeg=1};
    \node[badge, text=sourceblue!85!black, above=2pt of k6] {indeg=2};
    \node[badge, text=sourceblue!85!black, below=2pt of k7] {indeg=2};

    \draw[peeled] (k1) -- (k2);
    \draw[peeled] (k1) -- (k3);
    \draw[active] (k2) -- (k4);
    \draw[active] (k3) -- (k4);
    \draw[active] (k3) to[bend left=14] (k6);
    \draw[active] (k4) -- (k5);
    \draw[active] (k5) -- (k6);
    \draw[active] (k5) -- (k7);
    \draw[active] (k4) to[bend right=18] (k7);
  \end{scope}

  % Right panel: Cyclic residual core stuck
  \begin{scope}[xshift=8.2cm]
    \node[paneltitle] at (2.1, 2.35) {(b) Cyclic core: $Q=[\,]$ while $\abs{L}<n$};

    \node[vdone]  (c0) at (0.0, 0.6) {$s$};
    \node[vstuck] (c1) at (1.6, 0.6) {$u$};
    \node[vstuck] (c2) at (3.2, 1.45) {$v$};
    \node[vstuck] (c3) at (3.2, -0.25) {$w$};

    \node[badge, text=gray!70!black, above=2pt of c0] {out};
    \node[badge, text=trapred!85!black, below left=1pt and -2pt of c1] {indeg=1};
    \node[badge, text=trapred!85!black, above=2pt of c2] {indeg=1};
    \node[badge, text=trapred!85!black, below=2pt of c3] {indeg=1};

    \draw[peeled] (c0) -- (c1);
    \draw[cyc] (c1) -- (c2);
    \draw[cyc] (c2) -- (c3);
    \draw[cyc] (c3) -- (c1);
  \end{scope}
\end{tikzpicture}
\end{center}

\begin{labbox}
Re-run the trace but break ties by \emph{larger} vertex number, or use a stack (LIFO) instead of a
queue. You will get a different valid order. Confirm by hand that whatever legal choices you make, the
result still has every edge pointing forward. The freedom in tie-breaking is exactly the
non-uniqueness we characterize in Section~\ref{ch04:uniqueness}.
\end{labbox}

\begin{theorem}[Correctness of Kahn]\label{ch04:kahn-correct}
On input $\Gph$, Algorithm~\ref{ch04:kahn-algo} outputs a valid topological order if and only if
$\Gph$ is a DAG; otherwise it outputs fewer than $n$ vertices, and the omitted vertices are exactly
those that lie on a directed cycle or are reachable from one.
\end{theorem}

\begin{proof}
\textbf{Invariant.} At every moment, for each vertex $w$, $\mathrm{indeg}[w]$ equals the number of
in-edges $(u,w)$ whose tail $u$ has \emph{not} yet been appended to $L$. This holds at start (no
vertex output, so $\mathrm{indeg}$ is the true indegree) and is preserved: appending $u$ decrements
exactly the counts of its out-neighbors, which are precisely the edges leaving the just-output $u$.

\textbf{Order is legal.} A vertex enters $Q$ only when $\mathrm{indeg}=0$, i.e.\ when \emph{all} its
predecessors are already in $L$. So when $w$ is appended, every $u$ with $(u,w)\in E$ precedes it:
$\pi(u)<\pi(w)$ for all edges. Thus $L$, read as positions, is a topological order of the vertices it
contains.

\textbf{DAG $\Rightarrow$ all $n$ output.} If $\Gph$ is acyclic, suppose some vertices never enter
$Q$; let $S$ be that nonempty set. The induced subgraph $\Gph[S]$ is a DAG, so by
Lemma~\ref{ch04:source-exists} it has a vertex $s$ with no in-edge \emph{from within $S$}. Every in-edge
of $s$ thus comes from $V\setminus S$, all of which were output, so by the invariant $s$'s count
reached $0$ and $s$ would have been enqueued -- contradiction. Hence $S=\emptyset$.

\textbf{Cycle $\Rightarrow$ some vertex stuck.} If $\Gph$ has a cycle $C$, no vertex of $C$ can ever
reach indegree $0$: each has a predecessor on $C$ that itself never gets output (induction around the
cycle). So at least the vertices of $C$ stay in $S$, and $\abs L<n$. The stuck set is exactly the
vertices from which the algorithm cannot ``finish,'' which is the set of vertices lying on or
downstream of the residual cyclic core.
\end{proof}

\begin{complexitybox}
\textbf{Kahn's algorithm.} Computing indegrees scans all edges once: $\Oh(m)$. The seeding loop is
$\Oh(n)$. In the main loop each vertex is enqueued/dequeued exactly once ($\Oh(n)$ total) and each
edge $(u,w)$ is relaxed exactly once when $u$ is dequeued ($\Oh(m)$ total). Total time
$\Oh(n+m)$, extra space $\Oh(n)$ for the indegree array and queue. With a min-heap instead of a queue
you get the \emph{lexicographically smallest} topological order in $\Oh(n+m+ n\log n)$.
\end{complexitybox}

\begin{trapbox}
The most common bug is using $\abs L=n$ versus ``queue empty'' to decide success. The loop always
ends with an empty queue, so emptiness tells you nothing. You must compare the count of \emph{output}
vertices against $n$. A second trap: decrementing an indegree below zero because the input has
\emph{parallel edges} or a self-loop $(v,v)$. A self-loop makes $v$ depend on itself -- an
immediate one-vertex cycle -- and correctly leaves $v$ stuck; do not special-case it away. Parallel
edges are fine as long as you count each edge once when seeding and once when relaxing.
\end{trapbox}

\section{Method 2: DFS Finish Times}

Depth-first search gives a second route that needs no indegree bookkeeping and naturally reports the
cycle as a back edge. The idea is almost a slogan: \emph{a vertex finishes after all vertices it can
reach}. So if we order vertices by \emph{decreasing} finish time, every edge points from an
earlier-finishing-later vertex to a later-finishing-earlier one -- forward in the order.

\begin{intuitionbox}
Run DFS. When the recursion for $u$ returns, every successor reachable from $u$ has already returned
(its subtree of the search was fully explored before we backed out of $u$). ``Returned earlier'' will
become ``later in the topological order.'' So: push each vertex onto a stack at the \emph{moment it
finishes}, then read the stack top-to-bottom. The last vertex to finish sits on top and goes first.
\end{intuitionbox}

We use the standard DFS coloring: \textsc{white} (undiscovered), \textsc{gray} (on the recursion
stack, discovered but not finished), \textsc{black} (finished). A \emph{back edge} is an edge from a
vertex to a gray ancestor; its existence is equivalent to a cycle.

\begin{algorithm}[H]
\DontPrintSemicolon
\KwIn{Digraph $\Gph=(V,E)$.}
\KwOut{A topological order, or report ``cycle.''}
\ForEach{$v\in V$}{$\mathrm{color}[v]\gets\textsc{white}$}
$S\gets$ empty stack\; \quad $\mathit{cyclic}\gets \textbf{false}$\;
\ForEach{$v\in V$}{
  \If{$\mathrm{color}[v]=\textsc{white}$}{\textsc{Visit}$(v)$}
}
\eIf{$\mathit{cyclic}$}{\Return ``cycle''}{\Return $S$ read top-to-bottom}
\BlankLine
\SetKwProg{Fn}{Function}{:}{}
\Fn{\textsc{Visit}$(u)$}{
  $\mathrm{color}[u]\gets\textsc{gray}$\;
  \ForEach{$w$ with $(u,w)\in E$}{
    \uIf{$\mathrm{color}[w]=\textsc{gray}$}{$\mathit{cyclic}\gets\textbf{true}$ \tcp*{back edge: $w$ is an active ancestor}}
    \ElseIf{$\mathrm{color}[w]=\textsc{white}$}{\textsc{Visit}$(w)$}
  }
  $\mathrm{color}[u]\gets\textsc{black}$\;
  push $u$ onto $S$ \tcp*{$u$ finishes; it precedes nothing still on the stack}
}
\caption{\textsc{DFSTopoSort}$(\Gph)$}\label{ch04:dfs-algo}
\end{algorithm}

\begin{whybox}
Why ``decreasing finish time'' is exactly a topological order. Consider any edge $(u,v)$ at the
moment DFS explores it from a gray $u$. There are three cases for $v$'s color:
\begin{itemize}
\item \textsc{white}: $v$ is undiscovered, so DFS recurses into it; $v$ finishes \emph{inside} $u$'s
call, hence $f(v)<f(u)$.
\item \textsc{black}: $v$ already finished before $u$ explored the edge, so $f(v)<f(u)$.
\item \textsc{gray}: $v$ is an ancestor still on the stack -- this is a back edge and means a cycle
$v\rightsquigarrow u\to v$, so no topological order exists.
\end{itemize}
In the two acyclic cases $f(v)<f(u)$, i.e.\ $u$ finishes \emph{after} $v$. Ordering by decreasing
$f$ then puts $u$ before $v$: the edge points forward.
\end{whybox}

\begin{theorem}[DFS topological correctness]\label{ch04:dfs-correct}
Algorithm~\ref{ch04:dfs-algo} reports ``cycle'' iff $\Gph$ has a directed cycle; otherwise the stack,
read top to bottom, is a valid topological order.
\end{theorem}

\begin{proof}
First, \emph{back edge $\Leftrightarrow$ cycle}. If the algorithm flags a back edge $(u,w)$ with $w$
gray, then $w$ is an ancestor of $u$ in the DFS tree, so there is a tree path $w\rightsquigarrow u$,
and together with $(u,w)$ this is a directed cycle. Conversely, suppose $\Gph$ has a cycle; let $x$
be the first vertex of the cycle that DFS discovers. The other cycle vertices are white when $x$ is
discovered, so by the \emph{white-path theorem} (every white vertex reachable from $x$ via a white
path becomes a descendant of $x$) the cycle vertex $y$ with $(y,x)\in E$ becomes a descendant of $x$
and is gray when its edge $(y,x)$ to the still-gray ancestor $x$ is explored. That edge is flagged.

If no back edge is flagged, every explored edge fell in the white/black cases, giving $f(v)<f(u)$ for
every edge $(u,v)$ by the case analysis in the box above. Pushing on finish puts later-finishing
vertices nearer the top; reading top-to-bottom is decreasing $f$, so for every edge the tail (larger
$f$) precedes the head (smaller $f$). That is a topological order.
\end{proof}

\subsection*{Worked trace}

Run DFS on the same seven-task graph, visiting vertices and out-neighbors in increasing numeric
order. Discovery/finish times $(d/f)$ on a single global clock:

\begin{center}
\renewcommand{\arraystretch}{1.15}
\begin{tabular}{c c c c c c c c}
\toprule
vertex & $1$ & $2$ & $4$ & $5$ & $6$ & $7$ & $3$ \\
\midrule
$d$ & $1$ & $2$ & $3$ & $4$ & $5$ & $7$ & $12$ \\
$f$ & $14$ & $11$ & $10$ & $9$ & $6$ & $8$ & $13$ \\
\bottomrule
\end{tabular}
\end{center}

Walking it: start $1$ ($d{=}1$); go to $2$ ($d{=}2$); $2\to 4$ ($d{=}3$); $4\to 5$ ($d{=}4$);
$5\to 6$ ($d{=}5$, finishes $f{=}6$ -- no out-edges); $5\to 7$ ($d{=}7$, $7$ has no out-edge,
finishes $f{=}8$); $5$ finishes ($f{=}9$); $4\to 7$ is black, $4$ finishes ($f{=}10$); $2$ finishes
($f{=}11$); back at $1$, explore $1\to 3$ ($d{=}12$): $3\to 4$ black, $3\to 6$ black, $3$ finishes
($f{=}13$); finally $1$ finishes ($f{=}14$).

Push-on-finish stack (bottom $\to$ top): $6,7,5,4,2,3,1$. Read top-to-bottom:
$$1,\;3,\;2,\;4,\;5,\;7,\;6.$$
Every edge points forward (check $3\to 6$: $3$ sits at position 2, $6$ at position 7 -- forward;
$5\to 7$: positions 5 then 6 -- forward). Note this differs from Kahn's $1,2,3,4,5,6,7$: both are
valid, and the difference is exactly the ordering freedom.

The search built the DFS tree below (solid tree edges; dashed non-tree edges, all pointing to
already-\textsc{black} vertices, so none is a back edge and the graph is acyclic). Vertex $3$ is
discovered late, directly from the root $1$, only after $2$'s entire subtree has finished.

\begin{center}
\begin{tikzpicture}[>=Latex,
  v/.style={circle,draw=nodeblue,fill=blue!6,thick,minimum size=8.5mm,inner sep=0pt},
  tree/.style={draw=edgegray,thick,->},
  nontree/.style={draw=trapred,thick,->,dashed}]
  \node[v] (n1) at (0, 0) {$1$};
  \node[v] (n2) at (-1.2, -1.25) {$2$};
  \node[v] (n3) at (1.2, -1.25) {$3$};
  \node[v] (n4) at (-1.2, -2.5) {$4$};
  \node[v] (n5) at (-1.2, -3.75) {$5$};
  \node[v] (n6) at (-2.2, -5.0) {$6$};
  \node[v] (n7) at (-0.2, -5.0) {$7$};
  \draw[tree] (n1)--(n2); \draw[tree] (n2)--(n4); \draw[tree] (n4)--(n5);
  \draw[tree] (n5)--(n6); \draw[tree] (n5)--(n7); \draw[tree] (n1)--(n3);
  \draw[nontree] (n4) to[out=-25, in=70] (n7);
  \draw[nontree] (n3) to[bend right=18] (n4);
  \draw[nontree] (n3) to[out=-65, in=-30, looseness=1.45] (n6);
\end{tikzpicture}
\end{center}

\begin{trapbox}
Two classic DFS mistakes. (1) \emph{Appending on discovery} instead of on finish, then reading the
list forward -- that gives a \emph{pre}order, which is \emph{not} topological (a child can be
discovered before a sibling that must precede it). You must record finish order and reverse it.
(2) Treating a gray neighbor the same as a black one. A black neighbor is a harmless cross/forward
edge; a gray neighbor is a back edge and the cycle witness. Conflating them either misses cycles or
falsely reports them. For \emph{undirected} DFS the rule differs (the edge back to your parent is not
a cycle), but in a digraph every gray hit is a real back edge.
\end{trapbox}

\begin{complexitybox}
\textbf{DFS topological sort.} One DFS over the graph: each vertex colored once, each edge examined
once, $\Oh(n+m)$ time. Space $\Oh(n)$ for colors and the output stack, plus $\Oh(n)$ recursion depth
(replace with an explicit stack to avoid blowing the call stack on a 10-million-vertex chain). Same
asymptotics as Kahn; pick by taste and by what else you need (DFS also hands you the SCC machinery
and cycle witness ``for free'').
\end{complexitybox}

\begin{connectionbox}
The two methods are dual. Kahn peels \emph{sources} (indegree-$0$, front of the order); DFS-reverse
effectively peels \emph{sinks} (a finished vertex has all successors done, so it is a current sink of
the unfinished subgraph). Run DFS on the \emph{reverse} graph and read finish order forward and you
recover a source-first order. This same finish-time-reversal idea is the engine of Kosaraju's
strongly-connected-components algorithm in Chapter~5: there the condensation of the SCCs is itself a
DAG, and you topologically sort \emph{that}.
\end{connectionbox}

\section{Uniqueness of the Order}\label{ch04:uniqueness}

When is the topological order forced? Intuitively, you have a choice precisely when two vertices are
\emph{incomparable} -- neither can reach the other -- because then you may place either first. The
clean characterization is in terms of a \emph{Hamiltonian path}.

\begin{definition}[Hamiltonian path of a DAG]\label{ch04:ham-def}
A directed path that visits every vertex exactly once.
\end{definition}

\begin{theorem}[Uniqueness]\label{ch04:uniqueness-thm}
A DAG $\Gph$ has a unique topological order if and only if it has a Hamiltonian path, i.e.\ iff there
is an edge between each pair of consecutive vertices in the (necessarily then unique) order. Equivalently,
the order is unique iff at every step of Kahn's algorithm the queue holds exactly one vertex.
\end{theorem}

\begin{proof}
($\Leftarrow$) Suppose the order $v_1,\dots,v_n$ has $(v_i,v_{i+1})\in E$ for all $i$. Then for any
valid order $\sigma$, the edge $(v_i,v_{i+1})$ forces $v_i$ before $v_{i+1}$ in $\sigma$, so $\sigma$
must list them in the order $v_1,\dots,v_n$ -- identical. Hence unique.

($\Rightarrow$) Contrapositive: suppose no Hamiltonian path exists, i.e.\ in the order
$v_1,\dots,v_n$ produced by any run there is some $i$ with $(v_i,v_{i+1})\notin E$. We claim swapping
$v_i$ and $v_{i+1}$ yields another valid order, proving non-uniqueness. The swap only reverses the
relative order of $v_i,v_{i+1}$; it threatens validity only via an edge between them. There is no
edge $v_i\to v_{i+1}$ by assumption, and no edge $v_{i+1}\to v_i$ either (that would point backward in
the original valid order, impossible). All other edges keep their relative orientation. So the
swapped sequence is a second valid order.

The Kahn restatement: the queue having $\ge 2$ vertices at some step means two sources of the residual
graph are mutually incomparable at that moment, so either can be emitted next -- two orders. If the
queue is always a singleton, every choice is forced, and the unique consecutive emissions
$u_1,u_2,\dots$ each satisfy: when $u_{i+1}$ became the lone source, it was made ready by removing
$u_i$ (else it was ready earlier and would have coexisted in the queue), so $(u_i,u_{i+1})\in E$ --
a Hamiltonian path.
\end{proof}

\begin{intuitionbox}
The order is unique exactly when the DAG is a single ``chain with shortcuts'': a backbone
$v_1\to v_2\to\cdots\to v_n$ (the Hamiltonian path) plus optionally extra forward edges. Any missing
backbone link gives you a place to swap, hence a choice. So: \emph{unique order} $\equiv$ \emph{total
order already}, and a DAG with a genuine partial order (some incomparable pair) has at least two
topological orders. The number of topological orders equals the number of \emph{linear extensions} of
the partial order, a quantity that is famously $\#\mathrm{P}$-hard to count in general.
\end{intuitionbox}

The two small DAGs below contrast a unique topological order forced by the Hamiltonian chain $a\to b\to c$ (left) against a branching DAG where incomparable vertices $b$ and $c$ allow three distinct linear extensions (right).

\begin{center}
\begin{tikzpicture}[>=Latex,
  v/.style={circle,draw=nodeblue,fill=blue!6,thick,minimum size=7.5mm,inner sep=0pt}]
  \begin{scope}
    \node[v] (a) at (0,0) {$a$};
    \node[v] (b) at (1.8,0) {$b$};
    \node[v] (c) at (3.6,0) {$c$};
    \draw[->,thick,edgegray] (a)--(b); \draw[->,thick,edgegray] (b)--(c);
    \draw[->,thick,edgegray] (a) to[bend left=40] (c);
    \node[structurepurple, font=\small] at (1.8, -1.65) {unique: $a,b,c$};
  \end{scope}
  \begin{scope}[xshift=62mm]
    \node[v] (x) at (0,0) {$a$};
    \node[v] (y) at (1.8,0.65) {$b$};
    \node[v] (z) at (1.8,-0.65) {$c$};
    \node[v] (w) at (3.6,0.65) {$d$};
    \draw[->,thick,edgegray] (x)--(y); \draw[->,thick,edgegray] (x)--(z);
    \draw[->,thick,edgegray] (y)--(w);
    \node[structurepurple, align=center, font=\small] at (1.8, -1.65) {not unique: $b,c$ swap\\ $a,b,c,d$ \emph{or} $a,c,b,d$ \emph{or} $a,b,d,c$};
  \end{scope}
\end{tikzpicture}
\end{center}

\section{Dynamic Programming on a DAG}

Here is the payoff. A topological order is the perfect evaluation order for any quantity that flows
along edges: process vertices left to right, and by the time you reach $v$ every predecessor's value
is final. Acyclicity guarantees no circular dependency, so a single linear sweep computes recurrences
that would loop forever on a cyclic graph. This is dynamic programming where the DAG \emph{is} the
DP-state graph.

\begin{focusbox}
The template. Fix a topological order $v_1,\dots,v_n$. To compute $\mathrm{dp}[v]$ defined by a
recurrence over $v$'s in-neighbors (or out-neighbors), sweep in topological order (resp.\ reverse).
Each edge is relaxed once, so the whole DP costs $\Oh(n+m)$ -- linear, regardless of how many paths
the DAG hides.
\end{focusbox}

\subsection{Counting paths}

How many distinct directed paths run from a source $s$ to a target $t$? On a general graph this can be
exponential and tangled by cycles; on a DAG it is one sweep.

Let $\mathrm{paths}[v]$ be the number of distinct paths from $s$ to $v$. Then
$$\mathrm{paths}[s]=1,\qquad \mathrm{paths}[v]=\sum_{(u,v)\in E}\mathrm{paths}[u].$$
Processing vertices in topological order guarantees every $u$ with $(u,v)\in E$ comes before $v$, so
its count is ready. (A path to $v$ is a path to some predecessor $u$ extended by the edge $u\to v$,
and distinct predecessors or distinct paths to them give distinct $s$--$v$ paths; this partition is
exactly the sum.)

\begin{algorithm}[H]
\DontPrintSemicolon
\KwIn{DAG $\Gph$, source $s$, target $t$, topological order $v_1,\dots,v_n$.}
\KwOut{Number of $s\to t$ paths.}
\ForEach{$v\in V$}{$\mathrm{paths}[v]\gets 0$}
$\mathrm{paths}[s]\gets 1$\;
\For{$i\gets 1$ \KwTo $n$ (in topological order)}{
  $u\gets v_i$\;
  \ForEach{$w$ with $(u,w)\in E$}{
    $\mathrm{paths}[w]\gets \mathrm{paths}[w]+\mathrm{paths}[u]$\;
  }
}
\Return $\mathrm{paths}[t]$\;
\caption{\textsc{CountPaths}$(\Gph,s,t)$}\label{ch04:countpaths}
\end{algorithm}

On the seven-task graph with $s=1$, sweeping order $1,2,3,4,5,6,7$:
$\mathrm{paths}[1]=1$; $\to 2{:}1$, $3{:}1$; from $2$, $4{:}1$; from $3$, $4{:}2$, $6{:}1$; from $4$,
$5{:}2$, $7{:}2$; from $5$, $6{:}1{+}2{=}3$, $7{:}2{+}2{=}4$. So there are $3$ paths from $1$ to $6$
and $4$ from $1$ to $7$. (Trace $1\!\to\!6$ by hand: $1{,}3{,}6$; $1{,}2{,}4{,}5{,}6$;
$1{,}3{,}4{,}5{,}6$ -- three, as computed.)

The diagram below places the seven-task DAG into its linearized left-to-right order alongside the DP table $\mathrm{paths}[v]$ accumulated across the sweep.

\begin{center}
\begin{tikzpicture}[
    >=Stealth,
    vnode/.style={circle, draw=sourceblue, fill=sourceblue!10, thick, minimum size=7.5mm, font=\small\bfseries},
    dedge/.style={draw=edgegray, thick, ->},
    cell/.style={draw=gray!70, thick, rounded corners=2pt, minimum width=1.1cm, minimum height=0.72cm, font=\small\ttfamily},
    ruler/.style={font=\tiny\ttfamily, text=gray!80!black},
    paneltitle/.style={font=\footnotesize\scshape, text=gray!80!black}
]
  % Lane 1: Linearized DAG in topological order
  \node[paneltitle, anchor=west] at (-0.8, 2.7) {Linearized DAG (every edge points strictly left-to-right):};
  
  \node[vnode] (u1) at (0.0, 1.4) {$1$};
  \node[vnode] (u2) at (1.5, 1.4) {$2$};
  \node[vnode] (u3) at (3.0, 1.4) {$3$};
  \node[vnode] (u4) at (4.5, 1.4) {$4$};
  \node[vnode] (u5) at (6.0, 1.4) {$5$};
  \node[vnode] (u6) at (7.5, 1.4) {$6$};
  \node[vnode] (u7) at (9.0, 1.4) {$7$};

  % Direct edges
  \draw[dedge] (u1) -- (u2);
  \draw[dedge] (u1) to[bend left=30] (u3);
  \draw[dedge] (u2) to[bend right=30] (u4);
  \draw[dedge] (u3) -- (u4);
  \draw[dedge] (u3) to[bend left=35] (u6);
  \draw[dedge] (u4) -- (u5);
  \draw[dedge] (u4) to[bend right=42] (u7);
  \draw[dedge] (u5) -- (u6);
  \draw[dedge] (u5) to[bend right=30] (u7);

  % Lane 2: DP path count array
  \node[paneltitle, anchor=west] at (-0.8, 0.0) {DP table accumulation $\mathrm{paths}[v] = \sum_{(u,v)\in E} \mathrm{paths}[u]$:};

  \node[cell, fill=sourceblue!15, draw=sourceblue!80] (p1) at (0.0, -0.9) {$1$};
  \node[cell, fill=sourceblue!15, draw=sourceblue!80] (p2) at (1.5, -0.9) {$1$};
  \node[cell, fill=sourceblue!15, draw=sourceblue!80] (p3) at (3.0, -0.9) {$1$};
  \node[cell, fill=targetgreen!15, draw=targetgreen!80] (p4) at (4.5, -0.9) {$2$};
  \node[cell, fill=targetgreen!15, draw=targetgreen!80] (p5) at (6.0, -0.9) {$2$};
  \node[cell, fill=quotientorange!15, draw=quotientorange!80] (p6) at (7.5, -0.9) {$3$};
  \node[cell, fill=structurepurple!15, draw=structurepurple!80] (p7) at (9.0, -0.9) {$4$};

  % Ruler index above cells
  \foreach \v/\x in {1/0.0, 2/1.5, 3/3.0, 4/4.5, 5/6.0, 6/7.5, 7/9.0} {
    \node[ruler] at (\x, -0.35) {$\mathrm{paths}[\v]$};
  }

  % Annotations below key cells
  \node[font=\tiny, text=targetgreen!90!black, below=2pt of p4] {$1+1$};
  \node[font=\tiny, text=targetgreen!90!black, below=2pt of p5] {$\mathrm{from\ } 4$};
  \node[font=\tiny, text=quotientorange!90!black, below=2pt of p6] {$1+2$};
  \node[font=\tiny, text=structurepurple!90!black, below=2pt of p7] {$2+2$};
\end{tikzpicture}
\end{center}

\subsection{Longest path}

Longest path is NP-hard on general graphs (it encodes Hamiltonicity), yet on a DAG it is linear. The
recurrence for the longest path \emph{ending} at $v$ (number of edges, or summed weight $w$):
$$\mathrm{len}[v]=\max\Bigl(0,\ \max_{(u,v)\in E}\bigl(\mathrm{len}[u]+w(u,v)\bigr)\Bigr).$$

\begin{whybox}
Acyclicity is doing the heavy lifting. On a graph with a positive-weight cycle, ``longest path'' is
unbounded (loop forever) and the recurrence has no fixpoint. On a DAG the dependency
$\mathrm{len}[v]\leftarrow\mathrm{len}[u]$ only ever points backward in topological order, so the
system is triangular and a single sweep solves it. The exact same triangularity is why DAG
\emph{shortest} paths need no Dijkstra/Bellman-Ford -- one relaxation pass per vertex suffices, and
negative weights are no obstacle because there are no cycles to exploit.
\end{whybox}

\begin{algorithm}[H]
\DontPrintSemicolon
\KwIn{Weighted DAG $\Gph$ with weights $w$, topological order $v_1,\dots,v_n$.}
\KwOut{$\mathrm{dist}[\cdot]$ = longest-path (or shortest, swapping $\max/\min$) distances from sources.}
\ForEach{$v\in V$}{$\mathrm{dist}[v]\gets 0$ \tcp*{longest; use $+\infty$ everywhere but the source for single-source shortest}}
\For{$i\gets 1$ \KwTo $n$ (topological order)}{
  $u\gets v_i$\;
  \ForEach{$x$ with $(u,x)\in E$}{
    \If{$\mathrm{dist}[u]+w(u,x) > \mathrm{dist}[x]$ \tcp*{$<$ for shortest paths}}{
      $\mathrm{dist}[x]\gets \mathrm{dist}[u]+w(u,x)$\;
      $\mathrm{pred}[x]\gets u$\;
    }
  }
}
\caption{\textsc{DAGRelax}$(\Gph,w)$}\label{ch04:dagrelax}
\end{algorithm}

\begin{theorem}[DAG single-source shortest/longest paths]\label{ch04:dag-sssp}
Given a topological order, Algorithm~\ref{ch04:dagrelax} computes correct shortest- (or longest-) path
distances from the chosen source(s) in $\Oh(n+m)$, for \emph{arbitrary} edge weights (including
negatives), with no cycle-handling machinery.
\end{theorem}

\begin{proof}
We prove the shortest-path version; longest is symmetric. Claim: when the sweep reaches $v=v_i$,
$\mathrm{dist}[v]$ already equals the true shortest distance $\delta(s,v)$. Induct on $i$. A shortest
$s$--$v$ path's last edge is $(u,v)$ for some $u$ with $\pi(u)<\pi(v)$ (every edge points forward), so
$u=v_j$ with $j<i$. By induction $\mathrm{dist}[u]=\delta(s,u)$ was finalized when $u$ was swept, and
at that moment we relaxed $(u,v)$, setting $\mathrm{dist}[v]\le\delta(s,u)+w(u,v)=\delta(s,v)$.
Conversely $\mathrm{dist}[v]$ is only ever set to the length of some genuine $s$--$v$ path, so
$\mathrm{dist}[v]\ge\delta(s,v)$. Equality holds; and no later relaxation lowers it because all of
$v$'s in-edges came from already-swept vertices. Each edge is relaxed once: $\Oh(n+m)$.
\end{proof}

\begin{labbox}
Take the seven-task graph as \emph{unweighted} (each edge length $1$) and compute the longest path
ending at each vertex from source $1$, sweeping $1,2,3,4,5,6,7$:
$\mathrm{len}[1]{=}0$, $[2]{=}1$, $[3]{=}1$, $[4]{=}\max(1{+}1,1{+}1){=}2$, $[5]{=}3$,
$[6]{=}\max(1{+}1,\ 3{+}1){=}4$, $[7]{=}\max(2{+}1,\ 3{+}1){=}4$. The longest chain to $7$ has length
$4$: $1\to 2\to 4\to 5\to 7$. This length is the \emph{critical path} -- exactly the minimum makespan
if every task takes unit time and independent tasks run in parallel.
\end{labbox}

\subsection{Why one sweep beats Dijkstra here}

\begin{intuitionbox}
Dijkstra needs nonnegative weights because it finalizes a vertex's distance as soon as it is the
cheapest unsettled vertex, relying on the fact that no later detour can lower that cost -- an assumption that fails when an edge has negative weight. Bellman-Ford pays $\Oh(nm)$ to tolerate negatives because it does not know a safe
processing order. A DAG \emph{hands you} a safe order: topological. Once you process predecessors
before successors, a single relaxation per edge is provably final. The structure of the graph replaces
the priority queue.
\end{intuitionbox}

\begin{complexitybox}
\textbf{DP on DAGs.} Topological sort $\Oh(n+m)$, then one sweep $\Oh(n+m)$, total $\Oh(n+m)$ and
$\Oh(n)$ memo space for: counting paths, longest path / critical path, shortest path with any weights,
number of distinct longest/shortest paths, and most ``flow-along-edges'' recurrences. Compare:
general-graph longest path is NP-hard; general shortest path with negatives is Bellman-Ford's
$\Oh(nm)$. The DAG assumption is worth an entire complexity class.
\end{complexitybox}

\section{Applications: Scheduling and Build Systems}

\begin{applicationbox}
\textbf{Build systems.} \texttt{make}, Bazel, Ninja, and friends model files as vertices and ``$a$ is
a prerequisite of $b$'' as an edge $a\to b$. A topological sort yields a legal compile order; the
\emph{critical path} (longest weighted path, where weights are per-target build times) is the minimum
wall-clock time achievable with unlimited parallelism, and tells you which target to optimize first.
A reported cycle is the dreaded ``circular dependency'' error -- and the DFS method even hands the
user the offending back edge.
\end{applicationbox}

\begin{applicationbox}
\textbf{Project scheduling (PERT/CPM).} Tasks with durations and precedence constraints form a
weighted DAG. \textsc{DAGRelax} in \emph{longest}-path mode computes the earliest start time of every
task (longest path from the start node); a reverse sweep computes latest start times; tasks whose
earliest and latest starts coincide lie on the critical path and have zero \emph{slack} -- delaying
any of them delays the whole project. This is the original setting Kahn was solving in 1962.
\end{applicationbox}

\begin{applicationbox}
\textbf{Spreadsheets, package managers, dataflow.} A spreadsheet recalculates cells in topological
order of the formula-dependency DAG; a circular reference is a detected cycle. Package managers
(\texttt{apt}, \texttt{npm}, Cargo) install dependencies bottom-up in topological order. Compilers
order instructions and resolve symbol/type dependencies the same way. Whenever a problem requires doing each step only after
its prerequisites are complete, it defines a DAG and calls for a topological sort.
\end{applicationbox}

\begin{connectionbox}
Topological order is the gateway to the rest of this book. Strongly connected components
(Chapter~5) turn \emph{any} digraph into a DAG of its SCCs and then topologically sort that
condensation -- so even cyclic graphs get a meaningful linear order at the component level. The DAG
shortest-path sweep here is the base case that Dijkstra (Chapter~7) and the all-pairs algorithms
(Chapter~9) generalize to cyclic graphs.

In Chapter~5, we explore this condensation directly: we partition general directed graphs into
their strongly connected components, see how Kosaraju and Tarjan extract them in linear time,
and use the resulting DAG of components to solve 2-SAT and dependency problems.
\end{connectionbox}

\section{Exercises}

\begin{exercise}[Implement and compare]
Implement both Algorithm~\ref{ch04:kahn-algo} and Algorithm~\ref{ch04:dfs-algo}. On a randomly
generated DAG with $n=10^6$ vertices and $m=3\cdot 10^6$ edges, verify they each run in well under a
second and that both outputs are valid topological orders (every edge points forward). Replace the
recursive DFS with an explicit stack and confirm it no longer overflows on a path of $10^6$ vertices.
\end{exercise}

\begin{exercise}[Lexicographically smallest order]
Modify Kahn's algorithm to output the lexicographically smallest topological order by replacing the
FIFO queue with a min-heap. Prove it is correct (a greedy exchange argument) and give its running
time. Then argue why the DFS method cannot easily be coerced into the lexicographically smallest
order.
\end{exercise}

\begin{exercise}[Cycle witness]
Extend the DFS method to \emph{report the actual cycle} (the list of vertices) when it finds a back
edge $(u,w)$, not just the boolean. Hint: keep the current gray recursion stack; the cycle is its
suffix from $w$ to $u$ plus the edge. Do the same for Kahn: after termination, reconstruct a cycle
among the stuck vertices in $\Oh(n+m)$.
\end{exercise}

\begin{exercise}[Counting linear extensions]
Write a recursive routine that counts \emph{all} topological orders of a small DAG. On the right-hand
example of Section~\ref{ch04:uniqueness} confirm the count is $3$. Explain why your routine is
exponential in the worst case and look up why counting linear extensions is $\#\mathrm{P}$-complete.
\end{exercise}

\begin{exercise}[Unique-order test in linear time]
Give an $\Oh(n+m)$ algorithm that decides whether a DAG's topological order is unique, using
Theorem~\ref{ch04:uniqueness-thm}. (Run Kahn and check the queue never holds two vertices, or verify
the produced order has an edge between every consecutive pair.) Prove your test correct.
\end{exercise}

\begin{exercise}[Longest path, all weights]
Prove Theorem~\ref{ch04:dag-sssp} handles arbitrary (including negative) weights correctly for both
shortest and longest paths, and explain precisely where the proof would break if the graph had a
single directed cycle. Construct a 3-vertex example showing longest path becomes ill-defined when one
edge is reversed to form a cycle.
\end{exercise}

\begin{exercise}[Critical path and slack]
On a weighted DAG with a single source and single sink, compute for every vertex its earliest start
(forward longest-path sweep) and latest start (reverse sweep), and define slack as their difference.
Prove a vertex lies on \emph{some} critical (longest source--sink) path iff its slack is $0$. Apply it
to the seven-task graph with unit weights and list the critical-path vertices.
\end{exercise}

\begin{exercise}[Break it: where DFS preorder fails]
Exhibit a DAG on which reading the DFS \emph{pre}order (discovery order) forward is \emph{not} a
topological order, and pinpoint the edge that points backward. Then show that reversing the
\emph{postorder} (finish order) fixes it. Use a 4-vertex graph.
\end{exercise}

\begin{exercise}[Minimum path cover, glimpse of matching]
A \emph{minimum path cover} of a DAG is the fewest vertex-disjoint directed paths covering all
vertices. Show it equals $n$ minus the size of a maximum matching in a bipartite ``split'' graph
(each vertex $v$ becomes $v_{\text{out}}$ on the left and $v_{\text{in}}$ on the right, with edge
$v_{\text{out}}u_{\text{in}}$ for each $(v,u)\in E$). This previews the matching--flow connection of
Chapters~15--18; compute a minimum path cover of the seven-task graph by hand.
\end{exercise}

\begin{exercise}[Dynamic topological order]
Suppose edges are inserted into a DAG one at a time and you must maintain a topological order after
each insertion (rejecting any edge that would create a cycle). Describe an algorithm faster than
recomputing from scratch each time, and bound its total work over a sequence of $m$ insertions.
(Look up the Pearce--Kelly online topological ordering algorithm to check your bound.)
\end{exercise}
